\chapterimage{chapter_head_2.pdf}
\chapter{Math Foundations: Summations \& Recurrences}

\section{Theoretical Core \& Formulas}

\subsection{Summations}
The process of adding things together, often represented by the capital Greek letter Sigma ($\Sigma$). To perform Big-O analysis, we seek a \textbf{closed-form} mathematical expression of the summation.

\begin{theorem}[Core Summation Formulas]
The following closed-form expressions are essential for algorithm analysis:
\begin{itemize}
    \item \textbf{Constant Series:} $\sum_{i=1}^n c = c \cdot n$
    \item \textbf{Arithmetic Series (Sum of $i$):} $\sum_{i=1}^n i = \frac{n(n+1)}{2} = O(n^2)$
    \item \textbf{Sum of Squares (Sum of $i^2$):} $\sum_{i=1}^n i^2 = \frac{n(n+1)(2n+1)}{6} = O(n^3)$
    \item \textbf{Sum of Cubes:} $\sum_{i=1}^n i^3 = \frac{n^2(n+1)^2}{4} = O(n^4)$
    \item \textbf{Geometric Series:} $\sum_{i=0}^{n-1} a r^i = \frac{a(r^n - 1)}{r - 1}$
    \item \textbf{Infinite Geometric Series (where $|x| < 1$):} $\sum_{i=0}^\infty x^i = \frac{1}{1-x}$
\end{itemize}
\end{theorem}

\subsection{Recurrences}
A recurrence relation models the running time $T(n)$ of a recursive function by describing how to calculate a term based on previous terms in the sequence.

\begin{definition}[Divide and Conquer General Form]
$T(n) = aT(n/b) + D(n) + C(n)$
\begin{itemize}
    \item $a$ is the number of subproblems.
    \item $b$ is the factor by which the input size is divided.
    \item $D(n)$ is the divide time, and $C(n)$ is the combine time.
\end{itemize}
\end{definition}

\section{Techniques}

\subsection{Master Theorem}
\begin{theorem}[Master Theorem]
Used to quickly determine the Big-O running time of divide-and-conquer algorithms matching the form $T(n) = aT(n/b) + f(n)$. It compares the work done in splitting/combining $f(n)$ to the work done in the recursive calls $O(n^{\log_b a})$.
\end{theorem}

\subsection{Iteration Method (Substitution Method)}
A step-by-step technique to unroll a recurrence relation into a summation, which can then be reduced to a closed form.
\begin{enumerate}
    \item \textbf{Expand} the recurrence step by step for at least 3 iterations.
    \item \textbf{Notice the pattern} that emerges and write a generalized equation in terms of $k$.
    \item \textbf{Determine the base case} to find the value of $k$ that terminates the recurrence.
    \item \textbf{Derive a closed-form solution} by substituting $k$ back into the generalized equation and simplifying any summations.
\end{enumerate}

\section{Algorithmic Tracing: Iteration Method Example}

\begin{example}[Solving $T(n) = 2T(n-1) + 5$]
\textbf{Problem:} Solve $T(n) = 2T(n-1) + 5$, with base case $T(1) = 1$.

\textbf{Step 1: Expand (Unroll)}
\begin{itemize}
    \item $k=1: T(n) = 2T(n-1) + 5$
    \item $k=2: T(n) = 2[2T(n-2) + 5] + 5 = 2^2T(n-2) + 2(5) + 5$
    \item $k=3: T(n) = 2^2[2T(n-3) + 5] + 2(5) + 5 = 2^3T(n-3) + 2^2(5) + 2(5) + 5$
\end{itemize}

\textbf{Step 2: Generalize Pattern}
\begin{itemize}
    \item $T(n) = 2^k T(n-k) + \sum_{i=0}^{k-1} 5 \cdot 2^i$
\end{itemize}

\textbf{Step 3: Determine Base Case}
\begin{itemize}
    \item We hit the base case when $n - k = 1$.
    \item Solving for $k$, we get $k = n - 1$.
\end{itemize}

\textbf{Step 4: Substitute and Solve (Closed-Form)}
\begin{itemize}
    \item Substitute $k = n - 1$ into the generalized equation:
    $T(n) = 2^{n-1}T(1) + 5 \sum_{i=0}^{(n-1)-1} 2^i$
    \item Knowing $T(1) = 1$, we get:
    $T(n) = 2^{n-1}(1) + 5 \sum_{i=0}^{n-2} 2^i$
    \item Apply the Geometric Series formula to the summation:
    $T(n) = 2^{n-1} + 5 \left( \frac{2^{n-1} - 1}{2 - 1} \right)$
    \item Simplify:
    $T(n) = 2^{n-1} + 5(2^{n-1} - 1) = 6(2^{n-1}) - 5 = 3 \cdot 2^n - 5$
\end{itemize}
\textbf{Result:} $O(2^n)$
\end{example}

\section{Skills \& Pitfalls}

\begin{remark}
\textbf{Off-by-one errors in summation indices:} Known closed-form formulas usually assume the index starts at 1 (or 0 for geometric series). If a summation begins at a different value (e.g., $\sum_{i=10}^n i$), you must shift the summation or subtract the missing prefix before applying the formula.
\end{remark}

\begin{remark}
\textbf{Misapplying the Master Theorem cases:} Ensure the recurrence exactly matches the required form $T(n) = aT(n/b) + f(n)$ and carefully determine whether $f(n)$ is polynomially smaller, larger, or equal to $O(n^{\log_b a})$.
\end{remark}

\subsection{Must-Know Skills}
\begin{itemize}
    \item \textbf{Index Manipulation:} The ability to algebraically manipulate limits or split summations is critical.
    \item \textbf{Logarithm Properties:} When evaluating base cases for divide and conquer algorithms (e.g., $n/2^k = 1 \implies k = \log_2 n$), exponent and logarithm properties are heavily utilized.
\end{itemize}
